\chapter{Trabalhos relacionados}
\label{cap:trabalhos-relacionados}

Os trabalhos analisados abordam os métodos de cálculo de demanda, as alternativas de fornecimento de energia elétrica para edificações coletivas e a atualização dos requisitos normativos das distribuidoras.

\section{Estudos sobre métodos de cálculo de demanda}
\label{sec:trabalhos-demanda}

\citeonline[capítulos~5 e~6]{souza2020demanda} aplicou os critérios de dez concessionárias ou grupos de distribuição à mesma edificação fictícia. As demandas estimadas variaram de 30,45 a 95,48~quilowatts (kW). A escolha do método para os apartamentos e a parcela sobre a qual incidia a margem de segurança ajudaram a explicar a variação.

\citeonline[seções~5 e~6]{santo2023demanda} aplicaram os procedimentos da Enel Ceará e da Enel Rio de Janeiro a três configurações residenciais. As diferenças de demanda foram mais expressivas nas edificações maiores e repercutiram na seleção dos condutores do ramal de ligação.

\citeonline[p.~42 e~71]{silva2026demanda} comparou a demanda prevista, associada à proteção geral instalada, com a maior potência aparente trifásica média em intervalos de dez minutos. Nos 28 edifícios residenciais estudados na região de Curitiba e no litoral do Paraná, o valor previsto superou o medido. Essa comparação ajuda a interpretar o alcance de uma estimativa de projeto quando faltam medições do empreendimento.

\section{Estudos sobre projetos elétricos de múltiplas unidades consumidoras}
\label{sec:trabalhos-projetos}

\citeonline[capítulos~3 a~5]{bento2021fornecimento} dimensionou e orçou três alternativas de fornecimento para condomínios residenciais fictícios com percursos internos de curta, média e longa extensão, todos com demanda inferior a 300~quilovolts-amperes (kVA). O estudo comparou redes aéreas compactas em média tensão, redes subterrâneas em baixa tensão e redes subterrâneas em média tensão com transformação interna. A rede aérea compacta apresentou o menor custo de implantação nos três cenários. No condomínio de maior extensão, a rede subterrânea em média tensão foi mais econômica que a alternativa subterrânea em baixa tensão, pois a localização dos transformadores em pedestal próximos aos blocos encurtou os trechos de distribuição em baixa tensão.

\section{Estudos sobre a evolução dos requisitos normativos}
\label{sec:trabalhos-evolucao}

\citeonline[seção~4.1]{santiago2011regulamentacao} examinou como a Coelce adaptou seus documentos à Resolução Normativa nº~414/2010 da Agência Nacional de Energia Elétrica (ANEEL). Entre os documentos analisados estava a proposta de revisão da NT-003 sobre fornecimento a edificações coletivas. O estudo fornece um antecedente cearense de atualização dos critérios de atendimento, medição, proteção e aprovação de projetos.

\section{Síntese dos trabalhos relacionados e delimitação da contribuição do trabalho}
\label{sec:sintese-trabalhos}

Os cinco trabalhos analisados abordam aspectos complementares do fornecimento de energia elétrica a edificações coletivas, mas nenhum deles compara diretamente as especificações 0126 e 1580. A contribuição deste trabalho consiste em sistematizar as alterações entre esses documentos e avaliar seus impactos sobre um projeto elétrico aprovado, abrangendo não apenas a demanda, mas também as condições de atendimento, proteção, medição, aterramento e documentação. Essa delimitação refere-se ao conjunto de estudos analisados, sem pressupor ineditismo absoluto.

O Quadro~\ref{qua:trabalhos-relacionados} sintetiza as contribuições dos trabalhos analisados.

\begin{quadro}[!hbp]
  \centering
  \begin{Spacing}{1.}
    \caption{Síntese dos trabalhos relacionados}
    \label{qua:trabalhos-relacionados}
  \end{Spacing}
  {\small
  \renewcommand{\arraystretch}{1.2}
  \setlength{\tabcolsep}{4pt}
  \begin{tabular}{|>{\raggedright\arraybackslash}p{0.25\linewidth}|>{\raggedright\arraybackslash}p{0.34\linewidth}|>{\raggedright\arraybackslash}p{0.33\linewidth}|}
    \hline
    \textbf{Trabalho} & \textbf{Objeto e base da análise} & \textbf{Contribuição ao tema} \\\hline
    Souza (2020) & Edificação fictícia; normas de diferentes distribuidoras. & O método adotado altera a demanda atribuída à mesma edificação. \\\hline
    Santo et al. (2023) & Edificações residenciais; normas da Enel Ceará e Rio de Janeiro. & A diferença entre procedimentos repercute na seção dos condutores. \\\hline
    Silva (2026) & Medições em 28 edifícios residenciais. & As estimativas superaram as medições na amostra. \\\hline
    Bento (2021) & Três condomínios fictícios. & A extensão do percurso influencia o custo relativo das alternativas subterrâneas. \\\hline
    Santiago (2011) & Atualização dos documentos da Coelce. & A revisão da NT-003 ilustra a adaptação dos critérios da distribuidora à regulação. \\\hline
    Presente trabalho & Comparação 0126 e 1580; projeto aprovado do Edifício Palatium. & Avaliação dos impactos das alterações normativas. \\\hline
  \end{tabular}}
  \Fonte{Elaborado pelo autor com base nos trabalhos discutidos neste capítulo.}
\end{quadro}
